\subsection{A complete-coefficient splitting and its exact simplicial limit}

\begin{proposition}[A-L08: Proved; relative injectivity over complete coefficients]
\label{AL08}
Let $p$ be any prime and put $S_p=\Gamma_{\mathbb Z_p}(T)$.
In the category of integer-weight-graded $S_p$-modules with
weight-preserving maps, every module $M$ on which multiplication
by $p$ is injective satisfies
\[
 \operatorname{Ext}^i_{\mathrm{gr}S_p}(M,S_p)=0
 \qquad(i>0).
 \tag{AL.40}
\]
There is no finite-generation assumption on $M$.

For an admissible theory $A$ over $\mathbb Z_p$ with
$p$-torsion-free homogeneous terms, choose
$N\geq\max\{1,\lambda_A(0)\}$ and write
$I_N=\bigoplus_{w\geq N}S_p^w$.
For each simplex $n$, the coefficient inclusion has a
weight-preserving $S_p$-linear cochain retraction
\[
 I_N[0]\longrightarrow A_n^{\geq N,*}
       \xrightarrow{\ r_n\ }I_N[0].
 \tag{AL.41}
\]
No family of these retractions can commute with all faces
and degeneracies.  The retractions in \textup{(AL.41)} are
maps of cochain modules; no multiplicative assertion is made.
\end{proposition}

\begin{proof}
Write $D=S_p/pS_p=\Gamma_{\mathbb F_p}(T)$.
Nagpal--Snowden, Theorem~7.1(a), printed p.~39 with
proof on p.~41 of arXiv:1606.03431v2,
shows that $D$ is injective in the integer-graded module
category \cite{IC-nagpal-snowden}.  Its hypothesis holds:
for the classical divided-power algebra the defining
parameter is zero in $\mathbb F_p$ at every index $p^a$,
$a\geq1$.

Choose a graded free resolution $P_\bullet\to M$.
The two-term free resolution
$S_p\xrightarrow{p}S_p$ of $D$ computes
$\operatorname{Tor}^{S_p}_1(M,D)=M[p]=0$ and
$\operatorname{Tor}^{S_p}_i(M,D)=0$ for $i>1$.
Consequently $P_\bullet/pP_\bullet$ resolves $M/pM$
by graded free $D$-modules.  There is an actual
isomorphism of Hom complexes
\[
 \operatorname{Hom}_{\mathrm{gr}S_p}(P_\bullet,D)
 \cong
 \operatorname{Hom}_{\mathrm{gr}D}
             (P_\bullet/pP_\bullet,D).
 \tag{AL.42}
\]
Injectivity of $D$ therefore gives
$\operatorname{Ext}^{i}_{\mathrm{gr}S_p}(M,D)=0$
for every $i>0$.

Set $S_a=S_p/p^aS_p$, for $a\geq1$.
The exact sequence
\[
 0\longrightarrow D\xrightarrow{\ p^{a-1}\ }S_a
          \longrightarrow S_{a-1}\longrightarrow0
 \qquad(a>1)
\]
and induction give $\operatorname{Ext}^i(M,S_a)=0$
for every $i>0$.  The same long exact sequence gives
surjectivity of
$\operatorname{Hom}(M,S_a)\to\operatorname{Hom}(M,S_{a-1})$.

Here is a direct justification of the inverse-limit step.
Let $C_a=\operatorname{Hom}_{\mathrm{gr}S_p}(P_\bullet,S_a)$.
Its transition maps are surjective in every cochain degree,
because each $P_i$ is projective.  The inverse limit is taken
in graded modules, hence separately in each weight, and
$S_p=\varprojlim_a S_a$ in that category.  Hom preserves
limits, so
$\varprojlim_a C_a=\operatorname{Hom}(P_\bullet,S_p)$.
The following is a short exact sequence of cochain complexes:
\[
 0\longrightarrow\varprojlim_a C_a\longrightarrow
 \prod_{a\geq1}C_a
 \xrightarrow{\ 1-\mathrm{shift}\ }
 \prod_{a\geq1}C_a\longrightarrow0.
 \tag{AL.43}
\]
Surjectivity on the right is obtained recursively using the
surjective transition maps.  Products of abelian groups are
exact, so the products have cohomology
$\prod_a H^i(C_a)$.  These groups vanish for $i>0$.
In degree zero the map $1-\mathrm{shift}$ is surjective
again, because the transitions on
$H^0(C_a)=\operatorname{Hom}(M,S_a)$ are surjective.
The long exact sequence of \textup{(AL.43)} therefore gives
$H^i\operatorname{Hom}(P_\bullet,S_p)=0$ for every $i>0$.
This proves \textup{(AL.40)} without any bound on the
number of free generators or on their weights.

We next construct \textup{(AL.41)}.  Put
$V_n=A_n^{\geq N,0}$ and
$W_n=\operatorname{im}(d:V_n\to A_n^{\geq N,1})$.
The vertical coefficient condition in form degree zero
gives a short exact sequence of graded $S_p$-modules
\[
 0\longrightarrow I_N\longrightarrow V_n
       \xrightarrow{\ d\ }W_n\longrightarrow0.
 \tag{AL.44}
\]
The module $W_n$ is $p$-torsion free, being a submodule
of the specified $p$-torsion-free term.  Moreover it
is supported in weights at least $N$.  Thus
$\operatorname{Hom}_{\mathrm{gr}S_p}(W_n,S_p/I_N)=0$.
Apply Hom from $W_n$ to
$0\to I_N\to S_p\to S_p/I_N\to0$.
This zero Hom group and \textup{(AL.40)} give
$\operatorname{Ext}^1(W_n,I_N)=0$.
Consequently \textup{(AL.44)} splits $S_p$-linearly.
Take its retraction in degree zero and the zero maps
in positive cochain degrees.  Since the target has
zero positive degrees and the source has no negative
degrees, these maps form the required cochain map.

Finally suppose such $r_n$ commuted with faces and
degeneracies.  Choose
$w\geq\max\{N,\mu_A(0)\}$ and take form degree zero.
They would give retractions of simplicial modules
\[
 \underline{S_p^w}\longrightarrow A^{w,0}
             \longrightarrow\underline{S_p^w}.
\]
The composite is the identity.  Applying normalized
homology in degree zero would factor the identity of
$S_p^w\cong\mathbb Z_p$ through
$H_0(N A^{w,0})=0$, which is impossible.
This proves the stated failure of a simplicial family
by the actual coefficient maps and their homology.
\end{proof}

The proof yields the asserted coefficient retractions
in each simplex.  It does not assert a contraction of
the entire vertical complex.  In particular it does not
provide the prescribed higher-form tail primitives
excluded by A-L02, A-L03, or AT12.
